% Full version only: the receive case. Source: design/receive.md and spec
% 10.3 (0.11.0); sizes from spec/scripts/sizes.md on the fee model of
% spec/fee-analysis.md. The credit leaf and the receive flow are shared
% with the short version's voucher section (body-credit.tex,
% body-receive.tex), so their numbers are fixed in one place.
\section{Receiving over Lightning}\label{sec:receive}

A user with no on-chain funds but a fee coin receives a Lightning payment and wants a
balance it can leave with alone. An Ark client gets an out-of-round
output co-signed by the operator that expires with its round, or, below
the operator's floor, a ledger credit with no on-chain claim. The tree
takes the same payment with less trust, through a \emph{cumulative}
voucher (Appendix~\ref{app:claim}): since a redemption advances the
source leaf's $\mathsf{seq}$ and voids every other voucher against it,
repeated receives are one voucher that the issuer re-signs for the
running total against the same claim leaf. The recipient keeps the
latest and needs no contest window, since it never prefers a smaller
state. The covenant stops the issuer from redeeming more than the source
holds; it does not stop over-issuance or voiding (advancing
$\mathsf{seq}$, or redeeming an older, smaller voucher, before the
recipient redeems). Both are closed by the source leaf's shape.

\paragraph{The dedicated credit leaf.}
% Shared by both versions: the dedicated credit leaf. Source: spec
% 7.5.3 and 10.3 (0.11.0). The heading and label are the includer's:
% a run-in paragraph of the voucher section (short) or of the receive
% section (full-receive.tex), so the text starts without a blank line.
A dedicated credit leaf is a source leaf that the issuer owns and whose
exit key is one recipient's, in a tree whose taptree carries
\tr{CLAIM}$^{x}$, \tr{UPDATE}$^{x}$ and the exact-payout exit, carries
\tr{REWRITE}$_h^{x}$ in place of \tr{REWRITE}$_h$ (or no rewrite at all),
and omits \tr{SPLIT}. \tr{CLAIM}$^{x}$ is \tr{CLAIM} with one further
signature over the transaction, by the source leaf's exit key: 7 bytes
of script and 64 of witness in FULL, the only tier that builds it. The
recipient's own redemptions carry that signature at no cost, and the
issuer cannot carve the leaf at all, not even with a voucher it signs to
itself. Over-issuance then has no second victim, since every voucher
pays the same recipient, who can watch the leaf's amount. Voiding needs
the leaf to change, and every transition that could change it needs the
exit key: \tr{CLAIM}$^{x}$, \tr{UPDATE}$^{x}$ and \tr{REWRITE}$_h^{x}$
by rule, while the exact exit and an eviction pay the whole balance to
that key and so void nothing the recipient is owed. The balance can only grow or pay its recipient, whose remaining
exposure is the issuer's liveness: a provider that stops issuing.

Three details keep it that way. The recipient checks from the taptree
root $m$ that no plain \tr{REWRITE}$_h$ is present, since it needs only
the owner's and the operator's signatures, both the issuer's. The issuer
funds the leaf only with credit it has issued or is about to, since any
excess is the recipient's on every path the issuer can take alone, and
tops it up by an amount-only increase under \tr{REWRITE}$_h^{x}$, which
needs only the operator's signature. And the leaf has a finite lease
with the operator as evictor, not a $\NUMS$ evictor, which every carve
would inherit, blocking the sweep of whatever cohort the claim leaf lands
in. The credit is one-way: if the recipient could also spend from it, an
older, larger voucher would be a stale state needing a window to
supersede, as in a bidirectional channel. To spend, the recipient
redeems into a channel leaf (\S\ref{sec:channels}).


% Shared by both versions: receiving over Lightning into a credit leaf.
% Source: spec 10.3 (the settle-first order, audit B2) and
% design/receive.md; sizes from spec/scripts/sizes.md on the model of the
% evaluation section (no payout output for a carve or an update).

\paragraph{The flow.} A user receives a Lightning payment into a credit
leaf that the operator owns and that names the user's key as exit key,
in the FULL tier, the only one with \tr{CLAIM}$^{x}$. The operator issues
the invoice and draws its preimage, so the payer's receipt proves that
the operator was paid, not the user; a user who wants a receipt naming
it hands the operator preimages in advance, and can then receive offline
only as many payments as it has pre-shared. The payer's HTLC reaches the
operator, which settles it \emph{first}; only then does it sign the user
a voucher for the new running total, a plain cumulative voucher against
the credit leaf. Nothing reaches the chain per payment and the user need
not be online, but a voucher draws only what the credit leaf holds: the
operator pre-funds it, with capital it recovers only with the user's
consent, or tops it up as credit accrues, about $122\vB$ per credit leaf
in a cohort \tr{REWRITE}$_h^{x}$~\cite{layertreerepo}, which voids no
voucher. The order matters. A
hash-locked voucher issued before the HTLC settles, with the preimage
held by the recipient, breaks the rule that the party the HTLC pays draws
the preimage (\iffullversion spec~9.6.2\else\S\ref{sec:hashlock}\fi),
and on a credit leaf it cannot be repaired: if the
HTLC then fails, retiring the voucher means advancing the credit leaf's
$\mathsf{seq}$, which needs the recipient's exit key, so the voucher
stays redeemable for a payment that never arrived. Settling first leaves
the recipient exposed for one payment, the one the operator has settled
and not yet vouched for.

\paragraph{Settlement.} Normally the voucher is never redeemed. The
operator places a leaf of the recipient's own in its next cohort graft,
at under $3\vB$ per member in a 256-member \tr{JOIN}$_8$, and retires the
voucher by advancing the credit leaf's $\mathsf{seq}$ and lowering its
amount by what it has paid out: one \tr{UPDATE}$^{x}$ per credit leaf
($1{,}048\vB$), which the recipient co-signs, or one
\tr{REWRITE}$_h^{x}$ for a cohort of credit leaves, about $148\vB$ per
changed leaf at $h = 6$ plus $16\vB$ for each exit-key signature. If the
operator stops cooperating, the recipient redeems its latest voucher
alone with one \tr{CLAIM}$^{x}$: $1{,}635\vB$ at $D = 16$, with no
timelock, against 1{,}311 to $2{,}511\vB$ and a relative timelock to
unroll an Ark round tree of 256 to 65{,}536 members (\S\ref{sec:eval}).
A carve conserves the tree's total and the exact exit pays in full, so
both fees come from a confirmed coin of the recipient's: with the exit,
$2{,}655\vB$, or 26{,}550\,sat at 10\,sat/vB, below which the fallback
costs more than it recovers (\iffullversion Appendix~\ref{app:eval}\else Appendix~\ref{app:impl}\fi).


Where one operator leaf backs many recipients instead, over-issuance is
equivocation with a compact proof, which a slashable bond can deter;
voiding has no compact proof, so no bond reaches it.
